\chapter{Complexity — U Halt Machine}
%%% generated by the blueprint pipeline from TCSlib.Complexity.CircuitComplexity.UHaltMachine
\section{Overview}

Declarations of the Lean module \texttt{TCSlib.Complexity.CircuitComplexity.UHaltMachine}.

\section{Declarations}

\begin{definition}[Binary expansion of a natural number]
\lean{Complexity.binaryExpansion}
\label{Complexity.binaryExpansion}
\leanok
For a natural number $n$, its binary expansion is the list of its binary digits in
most-significant-first order.
\end{definition}

\begin{definition}[Natural number encoded by a bit string]
\lean{Complexity.numOfString}
\label{Complexity.numOfString}
\leanok
For a finite Boolean string $s$, define its encoded natural number by starting with $1$
and reading the bits of $s$ from right to left, replacing the current value $m$ by
$2m+1$ for a true bit and by $2m$ for a false bit.
\end{definition}

\begin{lemma}[Positivity of a Boolean bit fold]
\lean{Complexity.foldr_bit_one_pos}
\label{Complexity.foldr_bit_one_pos}
\leanok
\difficulty{4}
For every finite list $xs$ of Boolean values, folding the natural-number bit operation
over $xs$ from right to left, starting with $1$, yields a positive natural number.
\end{lemma}

\begin{lemma}[Bits of a folded binary word]
\lean{Complexity.bits_foldr_bit_one}
\label{Complexity.bits_foldr_bit_one}
\leanok
\uses{Complexity.foldr_bit_one_pos}
\difficulty{4}
For every finite bit string $xs$, fold its bits from right to left using $b,n\mapsto
2n+b$, starting from $1$. The binary expansion of the resulting natural number, written
from least significant bit to most significant bit, is $xs$ followed by a final bit $1$.
\end{lemma}

\begin{theorem}[Binary encoding of bit strings]
\lean{Complexity.binaryExpansion_numOfString}
\label{Complexity.binaryExpansion_numOfString}
\leanok
\uses{Complexity.binaryExpansion, Complexity.bits_foldr_bit_one, Complexity.numOfString}
\difficulty{3}
For every finite bit string $s$, form a natural number by folding its bits from right to
left using the binary operation $b,n\mapsto 2n+b$, starting from $1$. The binary
expansion of this number, written with the most significant bit first, is the bit $1$
followed by $s$.
\end{theorem}

\begin{lemma}[Primitive Recursion of the Replicated Encoding]
\lean{Complexity.primrec_replicate_numOfString}
\label{Complexity.primrec_replicate_numOfString}
\leanok
\uses{Complexity.numOfString}
\difficulty{5}
For a Boolean word $s$, let $m(s)$ be the natural number obtained by starting with $1$
and reading the bits of $s$ in reverse order, replacing each current value $v$ by $2v+b$
for the bit $b\in\{0,1\}$. The function sending $s$ to the list of $m(s)$ copies of true
is primitive recursive.
\end{lemma}

\begin{definition}[Universal halting language]
\lean{Complexity.UHALT}
\label{Complexity.UHALT}
\leanok
\uses{Complexity.HALT, Complexity.binaryExpansion, Language.unary, Turing.MachineCode}
Given a machine representation scheme, the universal halting language consists of unary
encodings of natural numbers $n$ for which the binary expansion of $n$ begins with $1$
and the remaining bit string describes an input on which the represented machine halts.
\end{definition}

\begin{theorem}[Unary halting language inclusion]
\lean{Complexity.UHALT_le_allOnes}
\label{Complexity.UHALT_le_allOnes}
\leanok
\uses{Complexity.UHALT, Language.allOnes, Language.unary_le_allOnes, Turing.MachineCode}
\difficulty{1}
For every machine representation scheme $c$, the unary halting language associated with
$c$ is contained in the language of all-ones bit strings.
\end{theorem}

\begin{theorem}[Unary representation of the halting predicate]
\lean{Complexity.replicate_numOfString_mem_UHALT_iff}
\label{Complexity.replicate_numOfString_mem_UHALT_iff}
\leanok
\uses{Complexity.HALT, Complexity.UHALT, Complexity.binaryExpansion_numOfString, Complexity.numOfString, Language.replicate_mem_unary_iff, Turing.MachineCode}
\difficulty{4}
Let $c$ be a machine representation and let $s$ be a finite bit string. Form the unary
word of length equal to the integer whose binary expansion is $1$ followed by $s$. This
word belongs to the unary language that records, at each such length, whether $c$
decodes some representation $\alpha$ for which the decoded machine halts on an input $x$
when given the paired string consisting of each bit of $x$ doubled, followed by $01$ and
then $\alpha$, if and only if there exist bit strings $\alpha$ and $x$ such that $s$ is that
paired string and the machine decoded from $c$'s input $\alpha$ halts on $x$.
\end{theorem}

\begin{theorem}[Universal halting has polynomial circuit size]
\lean{Complexity.UHALT_mem_PPoly}
\label{Complexity.UHALT_mem_PPoly}
\leanok
\uses{BoolCircuit.PPoly, Complexity.UHALT, Language.unary_inPPoly, Turing.MachineCode}
\difficulty{1}
For every machine representation scheme $c$, the unary language of numbers whose binary
expansions begin with $1s$, where $s$ encodes an advice string and an input on which the
machine decoded from that advice halts, has Boolean circuits of fan-in at most two and
size bounded by a polynomial in the input length.
\end{theorem}

\begin{theorem}[Computability of the halting function]
\lean{Complexity.HALT_computable_of_decidesInTime_UHALT}
\label{Complexity.HALT_computable_of_decidesInTime_UHALT}
\leanok
\uses{Complexity.Computable, Complexity.HALT, Complexity.UHALT, Complexity.numOfString, Complexity.primrec_replicate_numOfString, Complexity.replicate_numOfString_mem_UHALT_iff, Turing.FinTM, Turing.FinTM.ComputesInTime, Turing.FinTM.DecidesInTime, Turing.FinTM.exists_comp_partial, Turing.MachineCode, Turing.MultiTapeTM.indicator, Turing.codePrim_machine}
\difficulty{5}
Let $c$ be a machine representation scheme, $M$ a finite-state Boolean Turing machine,
and $T$ a function from natural numbers to natural numbers. If $M$ decides the unary
language associated with the halting function of $c$ within time bound $T$, then some
finite-state Boolean Turing machine computes the function sending each bit string $s$ to
the one-bit string whose bit is the value of that halting function at $s$.
\end{theorem}

\begin{theorem}[Time-bounded undecidability of unary halting]
\lean{Complexity.not_decidesInTime_UHALT}
\label{Complexity.not_decidesInTime_UHALT}
\leanok
\uses{Complexity.HALT_computable_of_decidesInTime_UHALT, Complexity.HALT_not_computable, Complexity.UHALT, Turing.EffectiveMachineCode, Turing.FinTM, Turing.FinTM.DecidesInTime}
\difficulty{1}
For every effective machine code $c$, every finite-state Turing machine $M$ over the
binary alphabet, and every function $T:\mathbb{N}\to\mathbb{N}$, $M$ does not decide the
unary language whose words have all bits equal to $1$ and whose length $n$ satisfies:
the binary expansion of $n$ is $1$ followed by a suffix $s$, and the halting predicate
for $c$ accepts $s$, within time bound $T$.
\end{theorem}

\begin{theorem}[Universal halting exceeds every time bound]
\lean{Complexity.UHALT_not_mem_DTIME}
\label{Complexity.UHALT_not_mem_DTIME}
\leanok
\uses{Complexity.DTIME, Complexity.UHALT, Complexity.not_decidesInTime_UHALT, Turing.EffectiveMachineCode}
\difficulty{3}
For every effective machine code scheme $c$ and every function
$T\colon\mathbb{N}\to\mathbb{N}$, the unary language consisting of all-one words of
length $n$ such that the most-significant-first binary expansion of $n$ is $1$ followed
by a string $s$ for which some advice string $\alpha$ and input $x$ satisfy
$s=\langle\alpha,x\rangle$ and the machine decoded from $\alpha$ halts on $x$ with some
output is not decidable by any finite-state binary Turing machine within time $C T(n)$
for any constant $C\in\mathbb{N}$.
\end{theorem}

\begin{theorem}[Universal halting is outside exponential time]
\lean{Complexity.UHALT_not_mem_EXP}
\label{Complexity.UHALT_not_mem_EXP}
\leanok
\uses{Complexity.EXP, Complexity.UHALT, Complexity.UHALT_not_mem_DTIME, Turing.EffectiveMachineCode}
\difficulty{3}
For every effective machine code scheme $c$, the unary language consisting of all-one
words of length $n$ such that the most-significant-first binary expansion of $n$ is $1$
followed by a string $s$ for which some advice string $\alpha$ and input $x$ satisfy
$s=\langle\alpha,x\rangle$ and the machine decoded from $\alpha$ halts on $x$ with some
output is not in $\mathrm{EXP}=\bigcup_{C\in\mathbb{N}}\mathrm{DTIME}(n\mapsto
2^{n^C})$.
\end{theorem}

\begin{theorem}[Universal halting is not in polynomial time]
\lean{Complexity.UHALT_not_mem_P}
\label{Complexity.UHALT_not_mem_P}
\leanok
\uses{Complexity.P, Complexity.P_subset_EXP, Complexity.UHALT, Complexity.UHALT_not_mem_EXP, Turing.EffectiveMachineCode}
\difficulty{1}
For every effective machine code scheme $c$, the unary language consisting of all-one
words of length $n$ such that the most-significant-first binary expansion of $n$ is $1$
followed by a string $s$ for which some advice string $\alpha$ and input $x$ satisfy
$s=\langle\alpha,x\rangle$ and the machine decoded from $\alpha$ halts on $x$ with some
output is not in $\mathrm{P}=\bigcup_{C\in\mathbb{N}}\mathrm{DTIME}(n\mapsto n^C+1)$.
\end{theorem}

\begin{theorem}[Polynomial circuits exceed exponential time]
\lean{Complexity.exists_mem_PPoly_not_mem_EXP}
\label{Complexity.exists_mem_PPoly_not_mem_EXP}
\leanok
\uses{BoolCircuit.PPoly, Complexity.EXP, Complexity.UHALT, Complexity.UHALT_mem_PPoly, Complexity.UHALT_not_mem_EXP, Turing.exists_effectiveMachineCode}
\difficulty{4}
There exists a language of binary strings that is decidable by a family of Boolean
circuits of polynomial size and is not decidable by any deterministic finite Turing
machine in time $c\,2^{n^d}$ for every choice of natural numbers $c,d$.
\end{theorem}

\begin{theorem}[Polynomial-size circuits exceed exponential time]
\lean{Complexity.PPoly_not_subset_EXP}
\label{Complexity.PPoly_not_subset_EXP}
\leanok
\uses{BoolCircuit.PPoly, Complexity.EXP, Complexity.exists_mem_PPoly_not_mem_EXP}
\difficulty{3}
There is a language of binary strings recognized by a family of Boolean circuits with
fan-in at most two and size at most $a(n+1)^k$ on inputs of length $n$, for some natural
numbers $a,k$, that is not decidable by any deterministic finite Turing machine in time
at most $c2^{n^d}$ for some natural numbers $c,d$.
\end{theorem}

\begin{theorem}[Polynomial-size circuits exceed polynomial time]
\lean{Complexity.PPoly_not_subset_P}
\label{Complexity.PPoly_not_subset_P}
\leanok
\uses{BoolCircuit.PPoly, Complexity.P, Complexity.PPoly_not_subset_EXP, Complexity.P_subset_EXP}
\difficulty{3}
There is a language of binary strings recognized by a family of Boolean circuits with
fan-in at most two and size at most $a(n+1)^k$ on inputs of length $n$, for some natural
numbers $a,k$, that is not decidable by any deterministic finite Turing machine in time
at most $c(n^d+1)$ for any natural numbers $c,d$.
\end{theorem}

\begin{theorem}[Strict circuit containment from inclusion]
\lean{Complexity.P_ssubset_PPoly_of_subset}
\label{Complexity.P_ssubset_PPoly_of_subset}
\leanok
\uses{BoolCircuit.PPoly, Complexity.P, Complexity.PPoly_not_subset_P}
\difficulty{3}
If every language in $\mathsf{P}$ is decidable by a polynomial-size Boolean circuit
family with fan-in at most two, then $\mathsf{P}$ is a proper subset of the class of all
such circuit-decidable languages.
\end{theorem}
